\begin{proposition}
\label{proposition-ML-characterization}
Let $M$ be an $R$-module.  Let $(M_i, f_{ij})$ be a directed system of finitely
presented $R$-modules, indexed by $I$, such that $M = \colim M_i$.  Let
$f_i:
M_i \to M$ be the canonical map.  The following are equivalent:
\begin{enumerate}
\item For every finitely presented $R$-module $P$ and module map $f: P
\to M$, there exists a finitely presented $R$-module $Q$ and a module
map $g: P \to Q$ such that $g$ and $f$ dominate each other, i.e.,
$\Ker(f \otimes_R \text{id}_N) = \Ker(g \otimes_R \text{id}_N)$
for every $R$-module $N$.
\item For each $i \in I$, there exists $j \geq i$ such that $f_{ij}: M_i
\to M_j$ dominates $f_i: M_i \to M$.
\item For each $i \in I$, there exists $j \geq i$ such that $f_{ij}: M_i
\to M_j$ factors through $f_{ik}: M_i \to M_k$ for all $k \geq
i$.
\item For every $R$-module $N$, the inverse system
$(\Hom_R(M_i, N), \Hom_R(f_{ij}, N))$ is Mittag-Leffler.
\item For $N = \prod_{s \in I} M_s$, the inverse system
$(\Hom_R(M_i, N), \Hom_R(f_{ij}, N))$ is Mittag-Leffler.
\end{enumerate}
\end{proposition}

\begin{proof}
First we prove the equivalence of (1) and (2).  Suppose (1) holds and let $i
\in I$. Corresponding to the map $f_i: M_i \to M$, we can choose $g:
M_i \to Q$ as in (1).  Since $M_i$ and $Q$ are of finite presentation,
so is $\Coker(g)$.  Then by Lemma \ref{lemma-domination}, $f_i : M_i
\to M$ factors through $g: M_i \to Q$, say $f_i = h \circ g$
for some $h: Q \to M$.  Then since $Q$ is finitely presented, $h$
factors through $M_j \to M$ for some $j \geq i$, say $h = f_j \circ h'$
for some $h': Q \to M_j$.  In total we have a commutative diagram
$$
\xymatrix{
  &  M  & \\
M_i \ar[dr]_g \ar[ur]^{f_i} \ar[rr]^{f_{ij}} &
& M_j \ar[ul]_{f_j} \\
  & Q  \ar[ur]_{h'} &
}
$$
Thus $f_{ij}$ dominates $g$.  But $g$ dominates $f_i$, so $f_{ij}$ dominates
$f_i$.

\medskip\noindent
Conversely, suppose (2) holds.  Let $P$ be of finite presentation and $f: P
\to M$ a module map.  Then $f$ factors through $f_i: M_i \to M$
for some $i \in I$, say $f = f_i \circ g'$ for some $g': P \to M_i$.
Choose by (2) a $j \geq i$ such that $f_{ij}$ dominates $f_i$.  We have a
commutative diagram
$$
\xymatrix{
P \ar[d]_{g'} \ar[r]^{f}           & M  \\
M_i \ar[ur]^{f_i} \ar[r]_{f_{ij}} & M_j \ar[u]_{f_j}
}
$$
From the diagram and the fact that $f_{ij}$ dominates $f_i$, we find that $f$
and $f_{ij} \circ g'$ dominate each other.  Hence taking $g = f_{ij} \circ g' :
P \to M_j$ works.

\medskip\noindent
Next we prove (2) is equivalent to (3).  Let $i \in I$.  It is always true that
$f_i$ dominates $f_{ik}$ for $k \geq i$, since $f_i$ factors through
$f_{ik}$.  If (2) holds, choose $j \geq i$ such that $f_{ij}$ dominates
$f_i$.  Then since domination is a transitive relation, $f_{ij}$ dominates
$f_{ik}$ for $k \geq i$. All $M_i$ are of finite presentation, so
$\Coker(f_{ik})$ is of finite presentation for $k \geq i$.  By Lemma
\ref{lemma-domination}, $f_{ij}$ factors through $f_{ik}$ for all $k \geq i$.
Thus (2) implies (3).  On the other hand, if (3) holds then for any $R$-module
$N$, $f_{ij} \otimes_R \text{id}_N$ factors through $f_{ik}
\otimes_R \text{id}_N$ for $k \geq i$.  So $\Ker(f_{ik} \otimes_R
\text{id}_N) \subset \Ker(f_{ij} \otimes_R \text{id}_N)$ for $k
\geq i$.  But $\Ker(f_i \otimes_R \text{id}_N: M_i \otimes_R N
\to M \otimes_R N)$ is the union of $\Ker(f_{ik} \otimes_R
\text{id}_N)$ for $k \geq i$.  Thus $\Ker(f_i \otimes_R
\text{id}_N) \subset \Ker(f_{ij} \otimes_R \text{id}_N)$ for
any $R$-module $N$, which by definition means $f_{ij}$ dominates $f_i$.

\medskip\noindent
It is trivial that (3) implies (4) implies (5).  We show (5) implies (3).  Let
$N = \prod_{s \in I} M_s$. If (5) holds, then given $i \in I$ choose $j \geq i$
such that
$$
\Im( \Hom(M_j, N) \to  \Hom(M_i, N)) =
\Im( \Hom(M_k, N) \to  \Hom(M_i, N))
$$
for all $k \geq j$.  Passing the product over $s \in I$ outside of the
$\Hom$'s
and looking at the maps on each component of the product, this says
$$
\Im( \Hom(M_j, M_s) \to  \Hom(M_i, M_s)) =
\Im( \Hom(M_k, M_s) \to   \Hom(M_i, M_s))
$$
for all $k \geq j$ and $s \in I$.  Taking $s = j$ we have
$$
\Im( \Hom(M_j, M_j) \to  \Hom(M_i, M_j)) =
\Im( \Hom(M_k, M_j) \to   \Hom(M_i, M_j))
$$
for all $k \geq j$.  Since $f_{ij}$ is the image of
$\text{id} \in  \Hom(M_j, M_j)$ under
$\Hom(M_j, M_j) \to  \Hom(M_i, M_j)$,
this shows  that for any $k \geq j$ there is $h \in \Hom(M_k, M_j)$
such that $f_{ij} = h \circ f_{ik}$.  If $j \geq k$ then we can take
$h = f_{kj}$. Hence (3) holds.
\end{proof}
